\section{向量加法：从顺序循环到线程网格}

\subsection{顺序基线与并行分解}

向量加法（vector addition）是数据并行程序中的“Hello, World!”。给定长度均为 $n$ 的向量 $A$ 与 $B$，目标是计算
\[
C[i]=A[i]+B[i],\qquad 0\le i<n.
\]
\par\noindent\begin{minipage}{\linewidth}
顺序 C++ 实现用一个循环依次完成 $n$ 次加法：

\begin{lstlisting}[style=cuda]
void vecAddCPU(const float* A, const float* B, float* C, int n) {
    for (int i = 0; i < n; ++i) {
        C[i] = A[i] + B[i];
    }
}
\end{lstlisting}
\end{minipage}\par

并行化时，不需要保留这个循环的顺序执行结构。把第 $i$ 次循环迭代定义为一个独立工作单元，再把工作单元 $i$ 分配给全局线程 $i$：
\[
\text{线程 }i\quad\longleftrightarrow\quad C[i]=A[i]+B[i].
\]
因此，原来的循环被整个线程网格替代。并行化前应先回答“每个线程究竟独立负责什么工作”，因为这一选择会同时决定线程数量、数据布局与后续优化空间。

\begin{keybox}[title={唯一写入者映射}]
输出元素 $C[i]$ 由唯一线程 $i$ 负责写入。不同线程不会同时修改同一输出位置，因此不需要原子操作或同步。
\end{keybox}

\subsection{CPU 与 GPU 的分离内存空间}

经典的离散 GPU 系统模型由彼此分离的主机与设备内存空间组成：CPU 连接\term{主机内存}{host memory}，GPU 连接\term{设备全局内存}{device global memory}。主机指针默认指向主机地址空间，设备指针指向设备地址空间；在基础 CUDA 编程模型中，二者不能被对方代码直接解引用。

\slidefig{0.94}{page-15.png}{设备内存管理的五个阶段}

完整数据流由五步构成：
\begin{enumerate}
  \item 在设备端分配空间；
  \item 把输入数据从主机复制到设备；
  \item 在设备端启动核函数并计算；
  \item 把结果从设备复制回主机；
  \item 释放设备端空间。
\end{enumerate}

即使系统提供\term{统一内存}{unified memory}或某些物理共享内存机制，数据仍需要在合适时机出现在 GPU 可高效访问的位置。显式分配与显式复制能够清楚呈现数据移动的成本与依赖关系。

\subsection{设备内存管理 API}

\subsubsection{分配与释放}

\begin{center}
\begin{tabularx}{0.96\textwidth}{>{\ttfamily}l X}
\toprule
\normalfont\textbf{CUDA API} & \textbf{作用与关键参数}\\
\midrule
cudaMalloc(void **devPtr, size\_t size) & 在设备全局内存中分配 \code{size} 字节，并把起始地址写入 \code{*devPtr}。返回值类型为 \code{cudaError\_t}。\\
\addlinespace
cudaFree(void *devPtr) & 释放由 \code{devPtr} 指向的设备内存区域。\\
\bottomrule
\end{tabularx}
\end{center}

对 $n$ 个 \code{float} 元素，字节数应写成
\[
S=n\cdot\code{sizeof(float)}.
\]
\par\noindent\begin{minipage}{\linewidth}
实际代码宜用 \code{size\_t} 保存 $S$，避免大规模数组下的整数溢出。

\begin{lstlisting}[style=cuda,numbers=none]
size_t bytes = static_cast<size_t>(n) * sizeof(float);
float *A_d = nullptr, *B_d = nullptr, *C_d = nullptr;

cudaMalloc(reinterpret_cast<void**>(&A_d), bytes);
cudaMalloc(reinterpret_cast<void**>(&B_d), bytes);
cudaMalloc(reinterpret_cast<void**>(&C_d), bytes);
// ... use the allocations ...
cudaFree(A_d);
cudaFree(B_d);
cudaFree(C_d);
\end{lstlisting}
\end{minipage}\par

\begin{warningbox}[title={设备指针的语义}]
\code{A\_d} 的数值是设备地址。主机代码可以把它传给 CUDA 运行时 API 或核函数，但不应在普通主机表达式中写 \code{A\_d[i]}。设备端数据访问应发生在核函数中。
\end{warningbox}

\subsubsection{主机-设备数据复制}

\par\noindent\begin{minipage}{\linewidth}
\code{cudaMemcpy} 的接口为
\begin{lstlisting}[style=cuda,numbers=none]
cudaError_t cudaMemcpy(void* dst, const void* src,
                       size_t count, cudaMemcpyKind kind);
\end{lstlisting}
\end{minipage}\par

第四个参数明确复制方向：

\begin{center}
\begin{tabular}{ll}
\toprule
\textbf{枚举常量} & \textbf{方向}\\
\midrule
\code{cudaMemcpyHostToHost} & 主机到主机（host to host）\\
\code{cudaMemcpyHostToDevice} & 主机到设备（host to device）\\
\code{cudaMemcpyDeviceToHost} & 设备到主机（device to host）\\
\code{cudaMemcpyDeviceToDevice} & 设备到设备（device to device）\\
\bottomrule
\end{tabular}
\end{center}

\par\noindent\begin{minipage}{\linewidth}
向量加法的输入与输出复制分别为
\begin{lstlisting}[style=cuda,numbers=none]
cudaMemcpy(A_d, A_h, bytes, cudaMemcpyHostToDevice);
cudaMemcpy(B_d, B_h, bytes, cudaMemcpyHostToDevice);
// kernel launch
cudaMemcpy(C_h, C_d, bytes, cudaMemcpyDeviceToHost);
\end{lstlisting}
\end{minipage}\par

读代码时可先检查 \code{dst} 与 \code{src} 是否和复制方向一致。例如，\code{cudaMemcpyDeviceToHost} 要求目标指针位于主机端、源指针位于设备端。
